\section{TonyPi 抱物穿行与状态机控制}
\subsection{总体方案与控制基础}
\subsubsection{方案演进与入口区分}
本部分的核心思路是保留人工录制且已经具备抱持效果的抓取动作，以动作组作为运动基本单元，将视觉观测作为动作选择与状态切换条件。机器人双足步态受负载、关节偏置及地面摩擦影响，固定动作次数容易累积航向误差，因此从开环演练逐步转向闭环控制。
\begin{table}[H]\centering\small
\caption{人形控制程序的用途与边界}
\begin{tabularx}{\linewidth}{p{36mm}YY}\toprule
入口 & 主要逻辑 & 本报告采用的地位\\\midrule
\code{open_loop.py} & stand、grab、定次数前进/左右转、release & 调试基线，不能补偿累计偏移\\
\code{competition.py} & 识别/抓取确认、固定或视觉路线、放置确认；支持人工stand到grab起步 & 综合反馈框架，需现场ROI与任务接口配置\\
\code{lane_trial.py} & 人工摆位后stand到grab，有限次数道路居中/纠偏 & 局部闭环试验，不包含完整两次路口任务\\
\code{maze_route/marker_route.py} & 抱起后依次靠近A、左转找B、右转找C、停止放置 & 最新独立路线实现；局部实机验证，完整路线待验收\\\bottomrule
\end{tabularx}\end{table}

不同入口具有不同的感知与任务假设，后续状态机以最新标记路线为分析对象。复用旧入口的动作、持物确认或心跳功能时，应明确完成移植，不能仅因文件同处仓库就视为已经集成。


\subsubsection{动作组格式与抱物步态组合}
\code{.d6a}文件实质是SQLite数据库，\code{ActionGroup}表按\code{Index}顺序保存动作帧。每帧由帧序号、持续时间（毫秒）及各\code{ServoN}位置刻度组成，历史文件存在16通道和18通道两种宽度。总线舵机刻度与头部PWM脉宽不是同一种单位，不应将“刻度500中位”套用到头部1200或1500的PWM参数。角度与位置的关系可写为
\begin{equation}
q_i=k_i(p_i-p_{0,i}),
\end{equation}
其中$p_i$为位置刻度，$p_{0,i}$为约定参考刻度，$k_i$包含方向和角度比例。刻度500只是常用中位参考，不能单凭此值认定机械关节角为零；本方案主要复用实录动作，不依赖完整关节绝对角标定。

令抓取动作最后一帧的上肢集合为$\mathcal A=\{6,7,8,14,15,16\}$，步态第$k$帧各舵机位置为$p_{k,i}^{\rm gait}$，则抱物组合动作定义为
\begin{equation}
p_{k,i}^{\rm carry}=\begin{cases}
p_{{\rm last},i}^{\rm grab},&i\in\mathcal A,\\
p_{k,i}^{\rm gait},&i\notin\mathcal A.
\end{cases}
\end{equation}
这一组合保持下肢帧序列、持续时间与其余通道不变，只替换六个手臂通道，减少转向步态破坏夹持姿态的可能性。历史版本分别组合过\code{go_forward}、\code{go_forward_fast}及快速转向动作；最新标记路线使用\code{grab_final_v2}和\code{catch_ball_*}，不能跨版本混用抓取末帧。

最新动作副本中的六个抱持刻度依次为566、908、40、435、108、935。\code{mkcopies.py}从设备原始动作生成独立副本，并逐帧核对未修改的腿部、索引和时间。放置动作仅在第一帧下蹲时保持抱紧手臂，后续降下、松开及站起帧保留原动作。运行时若首帧与当前姿态差值超过25刻度，将过渡时长延长到300\,ms；这一过渡不等于动态稳定性已获证明。

针对机体前倾，仓库另存两组大腿舵机相反方向的20刻度偏置候选，均带\code{_offset}标记。该方案未作为当前机器人的默认步态，也没有完整实机稳定性结论，不能在报告中写成最终校正量。

\subsubsection{抓取、持物观察与头部姿态}
早期通用视觉框架采用HSV/LAB分割、形态学去噪及轮廓筛选；物块改为红色后使用跨越HSV色相两端的双阈值。通用框架可在抓取后检查持物ROI，失败时经复位再识别，并在放置后检查手中区域与地面区域。人工起步模式则从stand直接衔接grab，省去搜索与接近阶段。

最新\code{marker_route}默认原地播放\code{grab_final_v2}后锁定上肢，\textbf{并未接入上述颜色持物确认、自动失败重抓和放置结果确认链}。因此“抓取动作完成”与“物块已可靠抓住”应分别记录，不能将通用框架的功能当作最新入口已验证的能力。比赛要求掉落后回到掉落点重新抱起；当前路线程序也没有完成该恢复闭环。

头部参数按程序和目标分别标定：早期道路/持物调试使用PWM servo1=1200，servo2保持当前值；最新标记路线默认pitch=1500、pan=1482，用来观察前方墙面标记。不同俯仰角会改变地面可见区域、物块遮挡与近距离码出画面的位置，报告应为每次运行登记真实参数，不能把两套姿态并写成同一配置。

\subsubsection{早期灰色挡板道路居中}
地面与挡板都为灰色，仅靠颜色阈值难以区分。早期道路算法使用边缘、线段、挡板支撑立柱与底边的几何关系，排除孤立地缝和地面胶带。抱持物块所在区域单独遮罩，避免把红色物块边缘误判为道路边界。

在近、远两个图像截面分别得到左右边界$x_L(y),x_R(y)$，可定义归一化横向误差与方向误差
\begin{align}
e_c&=\frac{x_L(y_n)+x_R(y_n)}{2W}-\frac12,\\
e_h&=\frac{x_L(y_f)+x_R(y_f)-x_L(y_n)-x_R(y_n)}{2W}.
\end{align}
其中$W$为图像宽度。组合误差超过阈值时选择左/右纠偏动作，误差处于容差内时执行一次前进，等待动作结束后重新观察。观测需连续稳定，不能在一张模糊帧上立即切换。挡板同色、遮挡、低对比和双足运动模糊仍可能导致无效道路判断；该方法不提供米制障碍距离或通用绕障规划。

\subsubsection{最新A/B/C固定标记识别与测距}
最新程序不依赖地面标识，而是在墙面布置三张已知模板标记。每张标记由$9\times9$格构成，外两格为黑边，中间$5\times5$格为数据，黑色正方形边长14\,cm。检测过程为自适应阈值、四边形轮廓筛选、透视展开、单元格采样，以及四个旋转方向的模板匹配；允许的数据位误差为1格。它们是程序自定义模板，\textbf{不能当作ArUco、AprilTag或TurtleBot所用ALVAR标签互换}。

以标记四角为三维已知点，结合相机内参$K$与畸变参数，通过PnP求取旋转$R$和平移$t$：
\begin{equation}
\lambda\begin{bmatrix}u\\v\\1\end{bmatrix}
=K\left[R\mid t\right]\begin{bmatrix}X\\Y\\Z\\1\end{bmatrix}.
\end{equation}
检测模块输出摄像头到标记中心的直线距离、方位角及墙面朝向。内参读取设备上的相机标定文件。码中心与镜头存在高度差时，PnP直线距离不能直接等同于沿地面水平距离，转角距离需用真实摆位复核。

\subsubsection{完整路线状态机与自适应纠偏}
路线状态依次为准备、抱起、走向A、左转寻找B、走向B、右转寻找C、走向C、停止与放置。距离阈值默认40\,cm；路口并非固定播放若干次就认定90度转弯，而是沿预定方向逐次转动，直到下一张目标码进入视野，再由直行阶段继续对中。每次转前后分别等待300\,ms，每次前进后等待200\,ms。

为降低步态大转角引起的左右摆动，程序对连续3帧观测取中位数，并根据近期单次转向后的方位变化估计转角$\hat\theta_{\rm step}$。纠偏门限采用
\begin{equation}
|\beta|>\max\!\left(8^\circ,0.6\hat\theta_{\rm step}\right),
\end{equation}
其中$\beta$为目标方位角。程序还从可见标记时的相邻步观测估计直行偏航，并在目标暂不可见时进行预测补偿。\textbf{标记不可见与相机无新帧是不同故障}：前者在计数约束下仍可能继续走，后者等待恢复。默认未见目标的盲行上限为25步，目标曾可见后连续丢失也有暂停阈值；这一实现与课程“丢失后原地重捕获”的严格要求存在差距，需在正式联调前收紧失效策略并验收。

当标记在近目标区域居中消失时，程序可使用最近观测判定已到达。这是处理近距离码从画面下缘离开的启发式，不能替代障碍测距。路口最多尝试12次转向；连续20步推进不足5\,cm、总步数过多或相机长时间无图像时触发暂停/退出。默认没有深度或测距输入，因此尚不能声称具备完整防碰撞保证。

\subsubsection{相机可靠性与可视化调试}
现场曾出现\code{/dev/video0}消失，内核提示USB枚举错误和线缆异常，重新插接后恢复。早期通用框架通过本地HTTP快照共享网页相机，避免重复打开已被占用的设备；最新\code{robust_camera.py}则按USB身份查找相机、处理网页相机释放和断线重连，只返回0.5秒内的新画面。两种相机管理模式应各自独立运行。

最新窗口显示运行阶段、目标标记框、标记ID、距离与方位等信息；支持S启动/恢复、P暂停、Q或Esc退出。每次运行保存\code{route.log}、\code{decisions.csv}及标注截图，用于核对状态转移的观测依据。截图中识别框和数值只证明该时刻的检测，不证明全程搬运成功。
比赛视频在184\,s处记录了相机识别窗口与路线日志。下面保留原始屏摄画面，用于核对视觉检测和转角搜索分支；完整视频取证说明见第\ref{sec:video-evidence}节。
\reportphoto{figures/humanoid/marker_debug.jpg}{比赛视频184\,s处的标记检测与路线日志（V184）}{原始视频截图。}
画面中的标记检测框与寻找下一码日志提供了现场调试证据，但屏摄反光和透视限制读数可靠性，因此本图不用于计算距离误差或动作次数。

\subsubsection{已有验证与后续验收}
最新仓库说明记录实机已跑通“抱起、走向A、在约39.6\,cm处开始左转”；完整A/B/C路线仍未实机验收\cite{maze}。仓库测试使用合成透视图验证标记识别与PnP，并用给定步长、偏航和转角模型演练路线；合成实验精度与路线完成不能作为实地成功率。

后续应重点记录：满载单步距离与偏航、左右转角分布、两次路口完成率、物块掉落/接触次数、相机断线频次及完整路线时长。抓取成功确认、丢码原地重捕获、掉落恢复、终点放置确认以及任务心跳适配，仍需在最新标记入口中形成一致闭环。
